\documentclass[a4paper,11pt]{article}
\usepackage[utf8]{inputenc}
\usepackage[french]{babel}
\usepackage[T1]{fontenc}
\usepackage{amsmath,amssymb,amsthm}
\usepackage{geometry}
\geometry{margin=2cm}
\usepackage{enumitem}
\usepackage{hyperref}
\usepackage{booktabs}
\usepackage{array}
\usepackage{xcolor}
\usepackage{listings}
\usepackage{titlesec}
\usepackage{fancyhdr}
\usepackage{lastpage}
\usepackage{graphicx}
\usepackage{etoolbox}
\AtBeginDocument{\shorthandoff{;:!?}}

\titleformat{\section}{\large\bfseries\color{blue!60!black}}{}{0pt}{}[\vspace{-0.5ex}\rule{\textwidth}{0.4pt}]
\titleformat{\subsection}{\bfseries\color{blue!40!black}}{}{0pt}{}
\titleformat{\subsubsection}{\itshape\color{blue!30!black}}{}{0pt}{}

\lstdefinestyle{monstyle}{
  frame=single,
  numbers=left,
  numbersep=5pt,
  breaklines=true,
  basicstyle=\small\ttfamily,
  backgroundcolor=\color{gray!5},
  rulecolor=\color{gray!40},
  columns=fullflexible,
  keepspaces=true,
  showstringspaces=false
}
\lstset{style=monstyle}

\pagestyle{fancy}
\fancyhf{}
\fancyhead[L]{\small STA211 -- Application MICE}
\fancyhead[R]{\small Fiche r\'ecapitulative}
\fancyfoot[C]{\thepage/\pageref{LastPage}}
\renewcommand{\headrulewidth}{0.4pt}

\theoremstyle{definition}
\newtheorem*{definition}{D\'efinition}
\theoremstyle{remark}
\newtheorem*{remark}{Remarque}
\newtheorem*{important}{\`A retenir}

\title{\textbf{Fiche r\'ecapitulative -- STA211}\\[4pt]
\large Application MICE~: imputation multiple sur les donn\'ees Diab\`ete}
\author{}
\date{10 juin 2026}

\begin{document}
\maketitle
\thispagestyle{empty}

\begin{abstract}
Cette fiche pr\'esente un cas pratique d'imputation multiple par
\'equations cha\^{i}n\'ees (MICE) sur le jeu de donn\'ees Diab\`ete
(Pima Indians Diabetes). Elle couvre l'analyse exploratoire du
dispositif de manquance, la mise en \oe uvre de l'algorithme MICE,
la v\'erification de la convergence, la comparaison des m\'ethodes
d'imputation, le pooling des mod\`eles, et l'analyse de sensibilit\'e.
\end{abstract}

\vfill
\tableofcontents
\newpage

% ===================================================================
\section{Introduction}
% ===================================================================

L'\'etude de cas applique la m\'ethode d'imputation multiple par
\'equations cha\^{i}n\'ees (MICE) aux donn\'ees \textbf{Diab\`ete}
(Pima Indians Diabetes Database). L'objectif est de traiter les
valeurs manquantes avant une classification supervis\'ee (pr\'ediction
du diab\`ete).

Donn\'ees~: \url{https://github.com/niharikagulati/diabetesprediction}

\subsection{Chargement des packages}

\begin{lstlisting}[language=R]
library(mice)
library(micemd)
library(FactoMineR)
library(parallel)
library(DescTools)
library(VIM)
library(mvtnorm)
library(funModeling)
\end{lstlisting}

\subsection{Importation}

\begin{lstlisting}[language=R]
load("diabetes.Rdata")
dim(diabetes)  # 768 x 9
\end{lstlisting}

Le jeu de donn\'ees contient 768 observations et 9 variables~:
8 pr\'edicteurs (Preg, Glucose, BP, Skin.Thick, Insulin, BMI, DPF,
Age) et la r\'eponse binaire Outcome.

% ===================================================================
\section{Analyse exploratoire}
% ===================================================================

\subsection{Liens entre variables}

\begin{lstlisting}[language=R]
summary(diabetes)
pairs(diabetes[, 1:8], cex = .2)
matcor <- cor(diabetes[,1:8], use = "pairwise.complete.obs")
PlotCorr(matcor)

boxplot(diabetes[,"Age"] ~ diabetes$Outcome)
sapply(colnames(diabetes)[1:8],
  FUN = function(xx, diabetes){
    boxplot(diabetes[,xx] ~ diabetes$Outcome,
            main = xx, ylab = xx, xlab = "outcome")
  }, diabetes = diabetes)
\end{lstlisting}

\subsection{Analyse du dispositif de manquance}

\textbf{Analyse univari\'ee}~: on utilise \texttt{aggr} du package
VIM pour visualiser la proportion de valeurs manquantes par variable.

\begin{lstlisting}[language=R]
res.aggr <- aggr(diabetes)
res.aggr$missings
cbind(res.aggr$tabcomb, res.aggr$percent)
\end{lstlisting}

\textbf{Analyse bivari\'ee}~: on utilise le V de Cramer entre les
indicatrices de manquance pour d\'etecter des structures.

\begin{lstlisting}[language=R]
var.na <- which(res.aggr$missing$Count > 0)
pattern <- is.na(diabetes[, var.na])
matcram <- PairApply(pattern, CramerV)
PlotCorr(matcram)
\end{lstlisting}

\textbf{Analyse multivari\'ee}~: une ACM sur le pattern de manquance.

\begin{lstlisting}[language=R]
MCA(pattern)
\end{lstlisting}

\subsection{Nature du m\'ecanisme de manquance}

On explore les liens entre valeurs manquantes et observ\'ees~:

\begin{lstlisting}[language=R]
# Une variable complete (Age) et une incomplete (Skin.Thick)
marginmatrix(diabetes[, c("Age", "Skin.Thick")])
# Rouge = individus avec valeur manquante sur Skin.Thick
# Les individus incomplets semblent plus ages

# Tous les couples
marginmatrix(diabetes[,-ncol(diabetes)], cex = .2, gap = 0)

# Matrixplot
matrixplot(diabetes, sortby = 8)
\end{lstlisting}

\begin{important}
  Si les individus avec valeur manquante sur une variable ont des
  distributions diff\'erentes sur une autre variable (compl\`ete),
  le m\'ecanisme est probablement MAR (Missing At Random).
\end{important}

\textbf{ACM sur donn\'ees discr\'etis\'ees}~:

\begin{lstlisting}[language=R]
d_bins <- discretize_get_bins(data = diabetes)
diabetes.cat <- discretize_df(diabetes, d_bins)
res.mca <- MCA(diabetes.cat, graph = FALSE,
               level.ventil = 0.04)
plot(res.mca, choix = "ind", invisible = "ind")
\end{lstlisting}

% ===================================================================
\section{Imputation multiple avec MICE}
% ===================================================================

\subsection{Param\`etres par d\'efaut}

La fonction \texttt{mice} du package \texttt{mice} utilise par
d\'efaut~:
\begin{itemize}
  \item $m = 5$ tableaux imput\'es.
  \item \texttt{maxit = 5} it\'erations.
  \item M\'ethodes par d\'efaut~: \texttt{pmm} (predictive mean
    matching) pour les variables num\'eriques, \texttt{logreg} pour
    les binaires, \texttt{polyreg} pour les multinomiales.
\end{itemize}

On lance MICE avec $m = 5$ et \texttt{maxit = 50}~:

\begin{lstlisting}[language=R]
res.mice <- mice(diabetes, maxit = 50, printFlag = FALSE)
tableau1 <- complete(res.mice, 1)
summary(tableau1)
\end{lstlisting}

\subsection{V\'erification de la convergence}

\begin{lstlisting}[language=R]
plot(res.mice)
\end{lstlisting}

Les courbes de convergence doivent montrer une stabilisation des
param\`etres au fil des it\'erations (pas de tendance monotone).

\subsection{Comparaison des distributions imput\'ees vs observ\'ees}

\begin{lstlisting}[language=R]
densityplot(res.mice)
\end{lstlisting}

Les distributions des valeurs imput\'ees (en rouge) doivent
ressembler \`a celles des valeurs observ\'ees (en bleu) si le
mod\`ele d'imputation est correct.

\subsection{Analyse par marginmatrix}

On compare les couples de variables avant et apr\`es imputation~:

\begin{lstlisting}[language=R]
diabetes.vim <- cbind.data.frame(complete(res.mice),
                                 is.na(diabetes))
colnames(diabetes.vim) <- c(colnames(diabetes),
  paste0(colnames(diabetes), "_imp"))

# Couple Age - Skin.Thick
par(mfrow = c(1, 2))
marginmatrix(diabetes[, c("Age", "Skin.Thick")])
marginmatrix(diabetes.vim[, c("Age", "Skin.Thick",
  "Age_imp", "Skin.Thick_imp")], delimiter = "_imp")

# Tous les couples
marginmatrix(diabetes.vim, delimiter = "_imp")
\end{lstlisting}

Les points imput\'es (en marron/jaune) doivent se situer dans le
nuage des points observ\'es.

\subsection{Overimputation}

L'overimputation consiste \`a traiter chaque valeur observ\'ee comme
manquante, l'imputer, et comparer l'imputation \`a la valeur r\'eelle.

\begin{lstlisting}[language=R]
nnodes <- detectCores() - 1

# Imputation avec m = 100 pour overimputation
res.mice.over <- mice.par(diabetes, m = 100,
  nnodes = nnodes, maxit = 20)

# Overimputation sur BMI (variable 6)
plotinds <- sample(seq(nrow(diabetes)), size = 30)
res.over <- overimpute(res.mice.over,
  plotinds = plotinds,
  plotvar = 6,
  nnodes = nnodes)
\end{lstlisting}

\begin{important}
  L'overimputation est un outil de diagnostic~: si les valeurs
  imput\'ees sont syst\'ematiquement \'eloign\'ees des valeurs
  observ\'ees, le mod\`ele d'imputation est inadapt\'e.
\end{important}

% ===================================================================
\section{Choix des m\'ethodes d'imputation}
% ===================================================================

\subsection{M\'ethodes disponibles pour MICE}

On examine les m\'ethodes par d\'efaut et on teste des alternatives
pour BMI~:

\begin{lstlisting}[language=R]
res.mice$method  # methodes par defaut
# On teste trois methodes pour BMI
par(mfrow = c(2, 2), mar = c(4, 4, 2, 1) + 0.1)
for (method.imp in c("rf", "norm", "norm.boot")){
  method <- res.mice$method
  method["BMI"] <- method.imp
  res.mice.tmp <- mice(diabetes, method = method,
    maxit = 20, printFlag = FALSE)
  print(densityplot(res.mice.tmp, ~ BMI, main = method.imp))
}
\end{lstlisting}

M\'ethodes test\'ees~:
\begin{itemize}
  \item \texttt{pmm} (d\'efaut)~: Predictive Mean Matching.
  \item \texttt{norm}~: r\'egression lin\'eaire bay\'esienne.
  \item \texttt{norm.boot}~: r\'egression lin\'eaire avec bootstrap.
  \item \texttt{rf}~: For\^et al\'eatoire.
\end{itemize}

\subsection{Comparaison par overimputation}

On compare \texttt{norm} et \texttt{rf} via overimputation~:

\begin{lstlisting}[language=R]
# norm
method <- res.mice$method
method["BMI"] <- "norm"
res.mice.over.norm <- mice.par(diabetes, m = 100,
  nnodes = nnodes, maxit = 20, method = method)
res.over.norm <- overimpute(res.mice.over.norm,
  plotinds = plotinds, plotvars = 6, nnodes = nnodes)

# rf
method["BMI"] <- "rf"
res.mice.over.rf <- mice.par(diabetes, m = 100,
  nnodes = nnodes, maxit = 20, method = method)
res.over.rf <- overimpute(res.mice.over.rf,
  plotinds = plotinds, plotvars = 6, nnodes = nnodes)
\end{lstlisting}

% ===================================================================
\section{Analyse apr\`es imputation}
% ===================================================================

\subsection{R\'egression logistique sur chaque tableau imput\'e}

On ajuste un mod\`ele de r\'egression logistique sur chaque tableau
imput\'e avec \texttt{with.mids}~:

\begin{lstlisting}[language=R]
fit <- with(res.mice,
  glm(Outcome ~ Preg + Glucose + BP + Skin.Thick
    + Insulin + BMI + DPF + Age, family = binomial))
length(fit$analyses)  # 5 modeles (un par tableau impute)
\end{lstlisting}

\subsection{Pooling des r\'esultats (r\`egles de Rubin)}

\begin{lstlisting}[language=R]
res.pool <- pool(fit)
summary(res.pool)
\end{lstlisting}

Le \texttt{pooling} combine les estimations des $m$ mod\`eles en
utilisant les r\`egles de Rubin~:
\begin{itemize}
  \item $\bar{\theta} = \frac{1}{m}\sum_{k=1}^m \hat{\theta}_k$
    (moyenne des estimations).
  \item Variance totale~: $T = W + (1 + \frac{1}{m})B$, o\`u $W$
    est la variance intra-imputation et $B$ la variance
    inter-imputation.
\end{itemize}

\subsection{Comparaison avec les cas-complets}

\begin{lstlisting}[language=R]
res.glm.cc <- glm(Outcome ~ ., family = binomial,
  data = diabetes)
summary(res.glm.cc)
\end{lstlisting}

On compare les coefficients et les \'ecarts-types entre l'analyse
sur cas-complets et l'analyse apr\`es imputation multiple. Des
diff\'erences importantes peuvent indiquer un biais de s\'election
d\^u aux donn\'ees manquantes.

% ===================================================================
\section{Analyse de sensibilit\'e}
% ===================================================================

L'analyse de sensibilit\'e par la m\'ethode du \textbf{delta}
(\emph{delta adjustment}) permet de tester la robustesse des
r\'esultats sous diff\'erents sc\'enarios pour le m\'ecanisme de
manquance (notamment l'hypoth\`ese MNAR).

On applique un d\'ecalage $\delta$ aux valeurs imput\'ees de
Skin.Thick~:

\begin{lstlisting}[language=R]
boxplot(diabetes$Skin.Thick)
delta <- c(0, 10, 20, 50)
imp.all <- vector("list", length(delta))
names(imp.all) <- delta
post <- res.mice.over$post
for (i in 1:length(delta)){
  d <- delta[i]
  cmd <- paste("imp[[j]][,i] <- imp[[j]][,i] +", d)
  post["Skin.Thick"] <- cmd
  imp <- mice(diabetes, post = post, maxit = 5,
              seed = i, print = FALSE)
  imp.all[[i]] <- imp
}
\end{lstlisting}

On compare les distributions imput\'ees pour $\delta = 0$ et
$\delta = 50$~:

\begin{lstlisting}[language=R]
par(mfrow = c(1, 2))
bwplot(imp.all[["0"]])
bwplot(imp.all[["50"]])
densityplot(imp.all[["0"]], lwd = 3)
densityplot(imp.all[["50"]], lwd = 3)
\end{lstlisting}

Si les conclusions de l'analyse changent peu malgr\'e un d\'ecalage
important, les r\'esultats sont robustes \`a l'hypoth\`ese MNAR.

% ===================================================================
\section{R\'esum\'e du pipeline}
% ===================================================================

\begin{enumerate}
  \item \textbf{Import}~: charger les donn\'ees et les packages
    (mice, VIM, micemd, etc.).
  \item \textbf{Analyse exploratoire}~: \'etudier les liens entre
    variables (corr\'elations, boxplots).
  \item \textbf{Analyse du pattern de manquance}~: univari\'ee
    (aggr), bivari\'ee (V de Cramer), multivari\'ee (ACM).
  \item \textbf{Nature du m\'ecanisme}~: marginmatrix, matrixplot.
  \item \textbf{Imputation multiple}~: lancer mice() avec les
    param\`etres choisis ($m$, maxit, m\'ethodes).
  \item \textbf{Diagnostics}~: convergence (plot), distribution
    (densityplot), marginmatrix, overimputation.
  \item \textbf{Analyse des tableaux imput\'es}~: with.mids + pool.
  \item \textbf{Comparaison avec cas-complets}~: v\'erifier la
    pr\'esence de biais.
  \item \textbf{Analyse de sensibilit\'e}~: delta adjustment pour
    tester l'hypoth\`ese MNAR.
\end{enumerate}

% ===================================================================
\section*{Questions cl\'es (QCM)}
% ===================================================================

\addcontentsline{toc}{section}{Questions cl\'es (QCM)}

\begin{enumerate}

  \item Quel package R impl\'emente l'imputation multiple par
    \'equations cha\^{i}n\'ees~?
    \begin{enumerate}
      \item VIM
      \item mice
      \item missMDA
      \item Amelia
    \end{enumerate}
    \textbf{R\'eponse~: b (package \texttt{mice}).}

  \item Quelle est la m\'ethode d'imputation par d\'efaut de
    \texttt{mice} pour les variables num\'eriques~?
    \begin{enumerate}
      \item r\'egression lin\'eaire bay\'esienne
      \item predictive mean matching (pmm)
      \item for\^et al\'eatoire
      \item moyenne conditionnelle
    \end{enumerate}
    \textbf{R\'eponse~: b (pmm).}

  \item Combien de tableaux imput\'es sont cr\'e\'es par d\'efaut
    par \texttt{mice()}~?
    \begin{enumerate}
      \item 3
      \item 5
      \item 10
      \item 20
    \end{enumerate}
    \textbf{R\'eponse~: b ($m = 5$ par d\'efaut).}

  \item \`A quoi sert la fonction \texttt{plot(res.mice)}~?
    \begin{enumerate}
      \item Visualiser la distribution des donn\'ees
      \item V\'erifier la convergence de l'algorithme MICE
      \item Afficher les valeurs manquantes
      \item Comparer les imputations
    \end{enumerate}
    \textbf{R\'eponse~: b (v\'erification de la convergence).}

  \item Quelle fonction permet de comparer les distributions des
    valeurs imput\'ees et observ\'ees~?
    \begin{enumerate}
      \item aggr
      \item densityplot
      \item marginmatrix
      \item both a and b
    \end{enumerate}
    \textbf{R\'eponse~: d (densityplot et marginmatrix).}

  \item Comment appelle-t-on la technique consistant \`a traiter
    chaque valeur observ\'ee comme manquante pour \'evaluer le
    mod\`ele d'imputation~?
    \begin{enumerate}
      \item Validation crois\'ee
      \item Bootstrap
      \item Overimputation
      \item Delta adjustment
    \end{enumerate}
    \textbf{R\'eponse~: c (overimputation).}

  \item Quelle fonction permet d'agr\'eger les r\'esultats de $m$
    mod\`eles ajust\'es sur les tableaux imput\'es~?
    \begin{enumerate}
      \item with.mids
      \item pool
      \item complete
      \item summary
    \end{enumerate}
    \textbf{R\'eponse~: b (pool).}

  \item Dans l'analyse de sensibilit\'e par delta adjustment, que
    repr\'esente $\delta$~?
    \begin{enumerate}
      \item Le nombre d'it\'erations suppl\'ementaires
      \item Le d\'ecalage ajout\'e aux valeurs imput\'ees
      \item Le seuil de d\'etection des atypiques
      \item Le taux de valeurs manquantes
    \end{enumerate}
    \textbf{R\'eponse~: b (d\'ecalage $\delta$ ajout\'e aux valeurs
    imput\'ees pour tester MNAR).}

  \item L'analyse de sensibilit\'e par delta adjustment permet de
    tester l'hypoth\`ese~:
    \begin{enumerate}
      \item MCAR
      \item MAR
      \item MNAR
      \item Aucune
    \end{enumerate}
    \textbf{R\'eponse~: c (MNAR).}

  \item Quel est l'avantage principal de l'imputation multiple par
    rapport \`a l'imputation simple~?
    \begin{enumerate}
      \item Elle est plus rapide
      \item Elle prend en compte l'incertitude due aux valeurs
        manquantes
      \item Elle ne n\'ecessite pas de mod\`ele
      \item Elle ne produit qu'un seul tableau
    \end{enumerate}
    \textbf{R\'eponse~: b (prise en compte de l'incertitude via la
    variance inter-imputation).}

\end{enumerate}

% ===================================================================
\begin{thebibliography}{9}
% ===================================================================

\bibitem{audigier2025}
  V.~Audigier.
  \newblock \emph{Application mice}.
  \newblock Cours CNAM, 2025.

\bibitem{vanbuuren2018}
  S.~van~Buuren.
  \newblock \emph{Flexible Imputation of Missing Data}.
  \newblock 2nd ed., CRC Press, 2018.

\bibitem{mice}
  S.~van~Buuren, K.~Groothuis-Oudshoorn.
  \newblock mice~: Multivariate Imputation by Chained Equations in R.
  \newblock \emph{Journal of Statistical Software}, 45(3), 2011.

\bibitem{micemd}
  V.~Audigier, F.~Husson.
  \newblock \emph{micemd~: Multiple Imputation by Chained Equations
    with Multilevel Data}.
  \newblock \url{https://CRAN.R-project.org/package=micemd}

\end{thebibliography}

\end{document}
